\slajdid{09_1-s044}
\begin{frame}[label=09_1-s044]
\gusttekst
\[\alert{1.}\quad k_{n,L}(V_{DD}-V_O)(V_{DD}+2\Delta-V_O)=k_{n,D}V_O(2V_I-2V_{Tn,D}-V_O)\]
Ako se u ovoj oblasti nalazi $V_{IH}$ diferenciranjem leve i desne strane i izjednačavanjem $\frac{\partial V_O}{\partial V_I}=-1$ dobijamo
\[k_{n,L}(V_{DD}+2\Delta-V_O)+k_{n,L}(V_{DD}-V_O)=-k_{n,D}(2V_I-2V_{Tn,D}-V_O)+k_{n,D}V_O(2+1)\]
\[k_{n,L}(2V_{DD}+2\Delta-2V_O)=k_{n,D}(2V_O-2V_I+2V_{Tn,D})\]
\[\alert{2.}\quad V_I=V_{Tn,D}+V_O-\frac{k_{n,L}}{k_{n,D}}(V_{DD}+\Delta-V_O)\]
Zamenom u početni izraz
\[k_{n,L}(V_{DD}-V_O)(V_{DD}+2\Delta-V_O)=k_{n,D}V_O\left(V_O-2\frac{k_{n,L}}{k_{n.D}}(V_{DD}+\Delta-V_O)\right)\]
\end{frame}
